\documentclass[12pt]{article}
\usepackage[utf8]{inputenc}
\usepackage[T1]{fontenc}
\usepackage{hyperref}
\title{Zigula Language Resources for K-12 Curriculum Development}
\author{Compiled by Claude Sonnet 5.5 (claude-sonnet-5-5) for Andr\'as Kornai}
\date{}
\begin{document}
\maketitle

\section{Introduction}
Zigula (ISO 639-3: ziw) is a Bantu language of Tanzania with about 380,000 speakers. It is not known to be taught in schools. A related variety, Mushungulu (code xma), is spoken in Somalia. Search results turned up little material designed for schools. This report lists only what the searches returned.

\section{Primary Educational Resources}
No primers, textbooks or school programmes for Zigula turned up in the searches. The Omniglot page below gives a brief orthographic overview that could help with early literacy design:
\begin{itemize}
\item Omniglot, Zigula (Chizigua): \url{https://www.omniglot.com/writing/zigula.htm}
\end{itemize}

\section{Linguistic and Reference Materials}
\begin{itemize}
\item Wikipedia, Zigula language: \url{https://en.wikipedia.org/wiki/Zigula_language}
\item Zigula overview from the DICE project at the University of Missouri (PDF): \url{https://dice.missouri.edu/assets/docs/niger-congo/Zigula.pdf}
\item Zigula-English Dictionary compiled by W. H. Kisbey for the Universities' Mission to Central Africa (London, 1906). Only an antiquarian bookseller listing was found, so the contents and any digital copy are unverified: \url{https://booksinbelgium.be/nl/b/zigula-english-dictionary-compiled-for-the-universities-mission-to-central-africa-kisbey-w-h-9591222}
\item Wikitongues language page for Zigula: \url{https://wikitongues.org/languages/ziw/}
\end{itemize}

\section{Religious and Community Resources}
One search snippet suggested a Zigula Scripture sample (Matthew 4) on bible.is under the code ZIWPBT. The page itself did not appear in the results, so it is not listed as a resource. No community organisation was found.

\section{Digital and Audio Resources}
\begin{itemize}
\item Zigula proverbs with Swahili, French and English translations (PDF, Tangaza University, 2019): \url{https://afriprov.tangaza.ac.ke/wp-content/uploads/2020/05/zigula_2019.pdf}
\end{itemize}

\section{Curriculum Development Context}
Zigula has no standard orthography. Teaching it in schools would first need an agreed writing system and new materials. The Kisbey dictionary, the proverb collection and the DICE overview could serve as starting points for vocabulary and texts. Swahili is the language of schooling in Tanzania.
\end{document}